% sections/10-outlook.tex
\section{What is not proved, and what remains}\label{sec:outlook}

\paragraph{The transfer principle.} Theorem~\ref{thm:transferintro} separates the two halves of the argument: the planar certificate,
a finite computation in the plane, and its passage to rectangular tori, which the theorem carries out for every planar
certificate of the form of~\cite{OAI:A-Fock-space-inequality-and-the-Laughlin-spectral-gap-September-24-2026}, with errors
bounded by explicit functions of the shorter period that decay at a Gaussian rate. It is the part of this paper on which
further work can build, in particular on the spinful hollow-core model discussed at the end of this section. For the
certificate of~\cite{OAI:A-Fock-space-inequality-and-the-Laughlin-spectral-gap-September-24-2026} the bounds are explicit:
evaluated numerically (Remark~\ref{rem:evaluation}), they give the thresholds $16.457$, $16.857$ and $18.685$ magnetic lengths
for a positive constant, a constant above $1/25$ and a constant within $10^{-5}$ of $\gs$ (Theorem~\ref{thm:openai}), and
the lower bound $\gl(L)$ tends to $\gs$ at the Gaussian rate $O(L^4\,2^{-L^2/20})$ of Remark~\ref{rem:rate}. The candidate proof of~\cite{HouWu} claims a uniform
bound of broader scope, on flat tori of every lattice shape and with every boundary twist, above an existential threshold;
the bounds here are restricted to rectangular tori and are explicit.

\paragraph{Oblique tori.} Theorem~\ref{thm:main} is proved for rectangular tori, uniformly in the aspect ratio; it says nothing
about tori with an oblique lattice of periods. The kernel theorem (Theorem~\ref{thm:kernel} and Proposition~\ref{prop:overfill})
holds for every modular parameter, but for non-rectangular $\tau$ it is a statement about the operators $B^{(\tau)}_n$ of
\eqref{eq:Bntau}, and we have not identified $\sum_n(B_n^{(\tau)})^*B_n^{(\tau)}$ with the periodized planar pair projector on an oblique
torus. The definition $\HT:=\LL_2(\Pp_2(P_1))$ makes sense for every lattice, and we expect most of the argument to carry over: the
proof of the Moyal identity (Theorem~\ref{thm:moyal}) uses only the irreducibility of the fine-lattice magnetic translations of
the centre-of-mass torus and the vanishing of the averaged image terms, and the image-sum calculus of Section~\ref{sec:square}
never uses the shape of the lattice. What would have to be redone is the part tied to the rectangular Landau gauge: the orbital
form (Proposition~\ref{prop:orbital}), the pair Gram matrix and its levels (Lemmas~\ref{lem:gram} and~\ref{lem:levels}), the
irreducibility argument of Lemma~\ref{lem:irred} written in an oblique orbital basis, and the lattice sums $\Ec_m$ and $\Xi_z$, whose
first shell can contain six vectors instead of four. A bound uniform in the modular parameter would allow the modular parameter to
be varied within the gapped phase, which is the setting in which the Hall viscosity appears as an adiabatic curvature over the
space of flat metrics~\cite{ASZ,Read09}; this is a separate problem, not addressed here.

\paragraph{Sharper constants.} The threshold is set almost entirely by the four-body term $2444\,\Ec_{22}(1+\Ec_{23})^2$ and, after
it, by the tail excess (Table~\ref{tab:constants}). The first comes from bounding the periodized four-body target $\mathcal T_4$ by
the pair energy through Cauchy--Schwarz over all $n_D$ triples of each block (Proposition~\ref{prop:parts}(d)); the second from the
crude bound $\abs{\bra m\Dis\ket m}\le\min\{1,\mathfrak g_m\}$ on the diagonal entries of the local Gram matrix. A sharper estimate of
$\mathcal T_4$, or the exact Laguerre values in $\Xi_z$, would lower the threshold; we have not attempted either. The bound on the
pair-level term is already far below the others, and so is the non-orthonormality term $\Cth\Ec_8(1+\Ec_8)$.

\paragraph{Finite tori below the threshold.} Below the threshold the method gives nothing. For instance at the square aspect ratio
and filling one third the bound is positive only for $N\ge15$; smaller systems are accessible to exact diagonalization, which is a
different kind of statement.

\paragraph{The spinful case.} For spin-$\frac12$ fermions with the hollow-core (pure $V_1$) interaction, at $\Nphi=2N$ flux quanta,
the zero modes are the Haldane--Rezayi states. We expect a spinful transfer principle of the form of Theorem~\ref{thm:transferintro}: a
planar certificate of the same form for the spinful model, with constant $\gamma>0$, should give a uniform gap above the zero
modes on rectangular tori. What it would need is the spinful version of each step: the spinful pair Gram matrix and its levels; the torus Moyal
identities for pair states with spin labels, which should follow from the same Schur argument, since the magnetic translations
act on the orbital factor only; and the three-body coefficient of the square of the spinful Hamiltonian, whose exact
periodization should follow from Lemma~\ref{lem:product}, which does not depend on internal labels, but whose planar coefficients
replace the $q_z$ by spin-dependent blocks that have to be computed. None of this is done here. In the contrapositive, such a
principle would show that no planar certificate of this form can exist for the spinful hollow-core model if a uniform gap above
its zero modes on rectangular tori is excluded.
